%\chapter{Beispiel für Formatierungen}

Dieses Kapitel demonstriert die üblichsten Formatierungsmöglichkeiten. Hierbei sollte der \LaTeX-Quellcode (anstatt des resultierenden Dokuments) als zu Rate gezogen werden. :-)


\verb|\textbf| \textbf{Ein formatierter Text} normaler Text \verb|\emph|  \emph{Ein formatierter Text} normaler Text \verb|\footnote| \footnote{Fussnote}. \verb|\enquote| \enquote{Anführungszeichen} oder >>Anführungszeichen<<

12~Byte, 6~kg, 100~EUR, 1--12, 299~792~458~m/s, 3~Liter, von \ldots bis \ldots

$12~Byte, 6~kg, 100~EUR, 299~792~458~m/s, 3~Liter$

\verb|12~Byte|, \verb|6~kg|, \verb|100~EUR|, \verb|1--12|, \verb|299~792~458~m/s|, \verb|3~Liter|, \verb|von \ldots bis \ldots|

Liste der recht­schreib­lich schwieri­gen Wörter\footnote{\url{https://www.duden.de/Liste-der-rechtschreiblich-schwierigen-Woerter}}.

Rechtschreibkontrolle - eine korrekte Rechtschreibung und Grammatik sind ein Muss!\footnote{\url{https://languagetool.org/de/}}.



\section{Aufzählungen}

\begin{itemize}
	\item a
	\item b
\end{itemize}

\begin{enumerate}
	\item eins 
	\item zwei
\end{enumerate}


\begin{description}
	\item[Beschreibung...] xyz 
	\item[Beschreibung...] zyx
\end{description}


\section{Gliederung -- Abschnitte, Unterabschnitte \& Absätze} \label{sec:structure}
Ein (Latex-)Dokument lässt je nach Dokumentenklasse (nicht jede Klasse unterstützt jede Untergliederung) unterteilen bzw. gliedern. In diesem Dokument stehen folgende Befehle zur Verfügung:
\begin{itemize}
	\item \verb|\chapter{...}|
	\item \verb|\section{...}      \label{sec:...}|
	\item \verb|\subsection{...}   \label{subsec:...}|
	\item \verb|\subsubsection{...}\label{subsubsec:...}|
	\item \verb|\paragraph{...}    \label{par:...}|
	\item \verb|\subparagraph{...} \label{subpar:...}|
\end{itemize}

\section{Referenzen}

\paragraph{Verweise (label + autoref)}
\verb|\autoref| \& \verb|\label| Text (\autoref{code:one}). Text (\autoref{fig:Chicken1}) Text \autoref{tab:tabneu} Text \autoref{sec:structure}.

\paragraph{Quellenangaben (cite)}
\verb|\cite| Text\cite{monk:2014:raspberry} Quelle: ~\cite{kofler:2015:raspberry}

\paragraph{Quellenangaben (textcite)}
\verb|\textcite| Text \textcite{monk:2014:raspberry} Quelle: ~\textcite{kofler:2015:raspberry}

\paragraph{Quellenangaben (footfullcite)}
\verb|\footfullcite| Text\footfullcite{monk:2014:raspberry} Quelle: ~\footfullcite{kofler:2015:raspberry}

\emph{Aplpe TV}\footnote{\url{http://www.aplpe.cmo}}


\section{Abbildungen}

Text (\autoref{fig:Chicken1}).
\begin{figure}[!hb]% hier: !hb
	\centering
	\includegraphics[width=0.4\linewidth]{images/logo}
	\caption{Chicken chien}\label{fig:Chicken1}%% anpassen
\end{figure}

Text (\autoref{fig:Chicken2} und \autoref{fig:Chicken1}).

\begin{figure}[!hb]% hier: !hb
	\centering
	\includegraphics[height=0.4\linewidth,angle=90]{images/logo}
	\caption{Bild 90 Grad drehen.}\label{fig:Chicken2}%% anpassen
\end{figure}


\section{Quelltext}
\verb|\lstinline| oder \verb|\verb|.

\paragraph{verb}
Bsp. \verb|int|, \verb|bool| (\autoref{code:one}).

\paragraph{lstlisting}

\lstset{language=C++}% C, TeX, Bash, Python
\begin{lstlisting}[%% anpassen
caption={Es ist eine alte Tradition, eine neue Programmiersprache mit einem Hello-World-Programm einzuweihen. Auch dieses Buch soll mit der Tradition nicht brechen, hier ist das Hello-World-Programm in C++}, label=code:one]
// Ein- und Ausgabebibliothek
#include <iostream>

int main(){// Hauptfunktion
	std::cout << "Hallo Welt!" << std::endl;// Ausgabe
	return 0;
}
\end{lstlisting}

\section{Tabellen neu}

(\autoref{tab:tabneu}).
\begin{table}[!ht]% hier: !ht
	\centering 
	\caption{Tabelle neu, gute Beschreibung einfügen}\label{tab:tabneu}%% anpassen
	\begin{tabular}{@{}rlc@{}}
	\toprule 
    \textbf{Nr.} & \textbf{Begriffe} & \textbf{Erklärung}\\
	\midrule
    1 & a1 & a2\\
    2 & b1 & b2\\
    3 & c1 & c2\\
    4 & a1 & a2\\
	\bottomrule
 	\end{tabular}
\end{table}


\section{Gleichungen}

$x$--$y$, \( x^2 + y^2 = 1 \)

(\autoref{eq:summe}).
\begin{equation}\label{eq:summe}%% anpassen
	\sum \limits_{i=1}^n i = \frac{n(n+1)}{2}
\end{equation}



